\scanpage{167}
\mspage{22}
plus généraux, groupes algébriques sur $k$. La construction en (B) de complexes de \unsure{jacobiennes|jacobiens|jacobienne} locales était, évidemment, \inserted{un peu} lié à cet espoir -- l'homologie donnerait peut-être ce que la cohomologie nous cèle ! Depuis lors, on dispose du formalisme de la cohomologie cristalline, et on sait \inserted{(Berthelot)} que (du moins \deleted{que} pour $X$ projectif et lisse) elle a les propriétés correctes. Si on l'utilise comme une sorte d'étalon sûr à quoi ``mesurer'' les autres cohomologies, on trouve ceci : la partie de la cohomologie cristalline $H^i_{\unsure{\mathrm{cris}|\mathrm{crs}}}(X)$ qui peut être décrite en termes de cohomologie fppf de $X$ à coefficients dans des $k$-groupes algébriques, n'est qu'une \deleted{faible} petite partie du $H^i$ ; de façon précise, utilisant la structure \inserted{très rigide} supplémentaire \unsure{de|des} $H^i$ (\unsure{module|modules} de t.f. sur \unsure{les' anneau|l'anneau} $W(k)$ des vecteurs de Witt) provenant de l'existence de l'hom de Frobenius $H^i\mathrel{\mathop{\unsure{\relbar\joinrel{\scriptscriptstyle\wedge}|\longrightarrow}}\limits^{F}}H^i$ (semi-linéaire) on trouve : que l'on reste toujours dans la partie ``de pente $\leq 1$'' (alors que les pentes possibles varient entre $0$ et $i$\ldots) Ce explique pourquoi pour $i=1$\unsure{.}
% archive sheet number: 167
